% Generated by scripts/assemble.py; corrections belong in source/corrections.json.
\examhead{113}{厦门大学 2024 年数学分析试卷}{111}
% p112-b01; source PDF page 112
\blockanchor{p112-b01}
\question{1}(15 分) 求极限 $\lim_{n\to \infty}\frac{1\cdot 1! + 2\cdot 2! + \cdots + n\cdot n!}{(n + 1)!}$.\par
% p112-b02; source PDF page 112
\par\Needspace{6\baselineskip}
\blockanchor{p112-b02}
\question{2}(20 分) 函数 $f(x)$ 在 $[0,1]$ 上有二阶连续导数, 且 $f(0) = f(1) = 0$, 当 $x \in (0,1)$ 时, $f(x) \neq 0$. 证明:\par
\nopagebreak[4]
% p112-b03; source PDF page 112
\blockanchor{p112-b03}
\displayformula{\int_ {0} ^ {1} \left| \frac {f ^ {\prime \prime} (x)}{f (x)} \right| \mathrm{d} x \geq 4.}
% p112-b04; source PDF page 112
\blockanchor{p112-b04}
\question{3}(15 分) 已知无穷积分 $\int_{a}^{+\infty} f(x) \, \mathrm{d}x$ 收敛, 且 $f(x)$ 单调, 证明: $\lim_{x \to +\infty} x f(x) = 0$.\par
% p112-b05; source PDF page 112
\blockanchor{p112-b05}
\question{4}(20 分) 求二元函数 $f(x, y) = x^2 + y^2 - 12x + 16y$ 在 $D = \{(x, y) \mid x^2 + y^2 \leq 25\}$ 上的最值.\par
% p112-b06; source PDF page 112
\blockanchor{p112-b06}
\question{5}(20 分) 有 $K$ 是 $n$ 维欧氏空间 $\mathbb{R}^n$ 上的有界非空闭子集, 且 $K \subseteq \bigcup_{k=1}^{\infty} U_k$, 其中 $U_k$ 为开集. 证明: 存在 $\varepsilon > 0$, 对于任意的 $x \in K$, 以 $x$ 为中心 $\varepsilon$ 为半径的开球 $B(x, \varepsilon)$ 包含于某个 $U_k$ 中.\par
% p112-b07; source PDF page 112
\blockanchor{p112-b07}
\question{6}(20 分) 解答如下问题:\par
% p112-b08; source PDF page 112
\blockanchor{p112-b08}
(1) 将 $f(x) = \begin{cases} -\frac{\pi}{4}, & x \in [-\pi, 0); \\ \frac{\pi}{4}, & x \in [0, \pi). \end{cases}$ 展开成傅里叶级数.\par
% p112-b09; source PDF page 112
\blockanchor{p112-b09}
(2) 求 $\sum_{n=1}^{\infty} \frac{1}{(2n-1)^2}$.\par
% p112-b10; source PDF page 112
\par\Needspace{6\baselineskip}
\blockanchor{p112-b10}
\question{7}(20 分) 设 $f(x)$ 有连续导数, 且 $f(0) = 0$, 求极限\par
\nopagebreak[4]
% p112-b11; source PDF page 112
\blockanchor{p112-b11}
\displayformula{\lim _ {t \rightarrow 0} \frac {1}{\pi t ^ {4}} \iiint_ {x ^ {2} + y ^ {2} + z ^ {2} \leq t ^ {2}} f \left(\sqrt {x ^ {2} + y ^ {2} + z ^ {2}}\right) d x d y d z.}
% p112-b12; source PDF page 112
\par\Needspace{6\baselineskip}
\blockanchor{p112-b12}
\question{8}(20 分) 计算曲线积分\par
\nopagebreak[4]
% p112-b13; source PDF page 112
\blockanchor{p112-b13}
\displayformula{\int_ {L} x \ln (x ^ {2} + y ^ {2} - 1) \mathrm{d} x + y \ln (x ^ {2} + y ^ {2} - 1) \mathrm{d} y.}
% p112-b14; source PDF page 112
\blockanchor{p112-b14}
其中 L 为从  $(2,0)$  到  $(0,2)$  的光滑曲线.\par
